\chapter{Introduction}
\label{ch:introduction}

\section{Motivation}

Geographic information systems (GIS) have become the standard instrument for
working with spatial data at every scale of public administration, science and
industry. The reason for their success is not only analytical power but
\emph{reproducibility}: an analysis performed in a desktop application can be
recorded as a script and re-executed on new data, on another machine, by
another analyst. For a long time this reproducibility was obtained from the
graphical user interfaces of systems such as GRASS \cite{neteler2012grass} and
QGIS \cite{qgissoftware}, whose Python bindings \cite{pyqgis,grass_tools} let
an analyst convert a click sequence into a re-runnable script after the fact.

That order of work, however, is inverted with respect to how analytical needs
are actually formulated. A practitioner usually knows the analysis in
\emph{intent} long before knowing its implementation: it is clear that the
streams must be buffered by 200 metres, or that average temperature must be
computed per administrative unit, while the exact module name, the spelling of
the parameters and the calling convention remain an implementation detail to be
looked up. The gap between intent and implementation is small in knowledge but
large in friction, and it is paid for on every single analysis.

Large language models (LLMs) appear to dissolve that gap: a model can be asked
for an analysis in ordinary language and can answer with code
\cite{vaswani2017attention,chen2021codex}. Since 2023 this capability has been
repackaged into \emph{agents}: systems in which a language model selects tools,
executes them, observes the outcome and iterates \cite{yao2023react}. The
geospatial community has adopted this direction vigorously, both through
research prototypes \cite{zhang2023geogpt,zhang2024geollmtooluse,singh2024geollmengine}
and through a rapidly growing ecosystem of QGIS plugins that put a chat window
in front of a running desktop session.

The enthusiasm, however, meets a hard constraint. The most capable code models
available in practice are large, and the most convenient way of running them is
as a hosted service. For an analyst handling sensitive spatial data --- land
cadastre, utility networks, health data, defence-related imagery --- neither the
upload of the prompt, nor the projection of the resulting code into a vendor's
cloud, nor the loss of control over which model version produced a given
analysis are acceptable. A large part of the current literature on LLM-assisted
GIS silently assumes a cloud endpoint \cite{mansourian2024chatgeoai,singh2024geollmengine}.
An offline alternative must therefore work with what a workstation can run: a
model of a few billion parameters, a local inference server, and the
documentation that ships with the software being automated. Offline-first
operation and hybrid routing are not mutually exclusive: keeping sensitive
requests local while allowing non-sensitive ones to reach a stronger hosted
model gives the analyst both control and capability.

This thesis addresses that gap. It develops and evaluates \textsc{GrassAgent},
an offline-first agent that translates unconstrained natural-language GIS
requests into executable Python scripts written against the modern
object-oriented \texttt{grass.tools} API of GRASS, verifies those scripts
deterministically against the official GRASS manual, and --- as a second stage
--- executes them in a live headless GRASS session and attempts to repair them
from the resulting error feedback.
